\begin{abstract}
El creciente número y complejidad de mega-constelaciones en órbita baja terrestre (LEO) exige operadores de misión más eficientes y resilientes. Este artículo explora el uso de Modelos de Lenguaje Grande (LLM) como interfaz de control intuitivo y agente de apoyo en operaciones de misión, facilitando la generación de comandos, diagnóstico de anomalías y toma de decisiones en lenguaje natural. Se propone una arquitectura agente LLM-aumentada que integra telemetría en tiempo real con razonamiento en lenguaje, evaluada en entornos de simulación de alta fidelidad. Los resultados demuestran mejoras significativas en tiempo de respuesta, reducción de errores humanos y usabilidad para operadores. Se discuten implicaciones para la autonomía en operaciones de constelaciones, desafíos de confiabilidad y caminos hacia implementación en vuelo, contribuyendo al avance de sistemas de operación de misión más intuitivos y seguros.
\end{abstract}